\subsection{Checking the projective threefold comparison}

In this subsection a \emph{good pair} is a smooth complex projective variety with reduced normal-crossings boundary, every connected open stratum, including $U$, aspherical, and every full normal-meridian lattice injecting into $\pi_1(U)$. The marking by $U$ is part of the data. The comparison below concerns these algebraically marked completions. It makes no assertion about arbitrary birational maps that move $U$.

\begin{result}{D.8}{Proved}{Good projective pairs are Mori hyperbolic and their comparisons are log crepant}
A good projective pair has no nonconstant algebraic map $\mathbb A^1\to S$ into any open stratum. For a simple normal-crossings boundary, $K_X+D$ is nef. For two such completions and a common smooth normal-crossings resolution $p:Z\to X$, $q:Z\to Y$, write $B=Z\setminus U$. Then
\[
 K_Z+B=p^*(K_X+D)+A_p=q^*(K_Y+E)+A_q,
 \qquad A_p=A_q\geq0.
\]
The common divisor is exceptional for both projections. A proper birational morphism between good smooth projective simple normal-crossings pairs, equal to the identity on $U$, is toroidal and log \'etale.
\end{result}
\begin{proof}
A map $f:\mathbb A^1\to S$ extends to $\bar f:\mathbb P^1\to X$ by properness. If infinity remains in $S$, asphericity makes the sphere nullhomotopic, whereas the integral of the pullback of a projective K\"ahler form is positive for a nonconstant holomorphic sphere. If infinity enters a deeper stratum $T$, trivialize the pulled-back old normal lines on the contractible source $\mathbb A^1$. Push the map a small nonzero distance in these directions into $U$, choosing the radius to tend to zero faster than every newly vanishing boundary coordinate near infinity. The boundary of a large disk then gives a nullhomotopic loop in $U$. In a normal-crossings chart at $\bar f(\infty)$ its newly acquired meridian coordinates are the nonzero contact orders of $\bar f$, with a common orientation sign. The old normal coordinates do not alter those new coordinates. This contradicts injection of the full lattice at $T$. Thus no such map exists. The open strata are exactly the strata in Svaldi's Theorem~1.1 \cite{Svaldi2019}; its hypotheses are smoothness, projectivity, reduced SNC boundary and the just-proved absence of these affine lines. Its conclusion gives nefness.

The logarithmic differential maps have effective determinant divisors $A_p,A_q$ supported on exceptional prime divisors. Put
\[
 R=p^*(K_X+D)-q^*(K_Y+E)=A_q-A_p.
\]
The divisor $-R$ is $p$-nef and $p_*R=p_*A_q$ is effective. The negativity lemma therefore gives $R\geq0$. Interchanging $p$ and $q$ gives $-R\geq0$, hence equality. This proves the claimed identity without asserting that the common resolution is good.

For an actual birational morphism $f:Y\to X$, the effective logarithmic Jacobian $A$ in $K_Y+E=f^*(K_X+D)+A$ is exceptional and $f$-nef. The same negativity lemma gives $A=0$. Thus the logarithmic differential matrix is everywhere invertible. In local boundary coordinates, $f^*x_i=u_i\prod_j y_j^{a_{ij}}$. The exponent matrix has full column rank; after choosing a nonsingular square minor and taking \'etale roots of units, absorb the corresponding units into the $y_j$. The remaining boundary-unit functions together with the ordinary target coordinates have independent ordinary differentials. They provide the invertible coordinates of a toric chart. Hence the morphism is locally an \'etale pullback of a monomial toric morphism; the finite lattice cokernel is invertible in characteristic zero. This is the explicit chart proof of log \'etaleness and toroidality.
\end{proof}

\begin{result}{D.9}{Proved}{Theorem R for projective threefolds, with the factorization reference corrected}
Two good smooth projective normal-crossings completions $(X,D)$ and $(Y,E)$ of the same smooth complex threefold $U$ admit a common good smooth projective completion $X_s$ such that $X_s\to X$ is a sequence of blow-ups of smooth closures of strata and $X_s\to Y$ is toroidal. They are joined over $U$ by a finite sequence of smooth stratal blow-ups and blow-downs through good projective pairs. This confirms Theorem~R of the attached review in its projective threefold scope.
\end{result}
\begin{proof}
First blow up the closures of the self-intersection strata in increasing dimension, obtaining simple normal-crossings boundaries. Locally these are barycentric subdivisions of the normal orthants. Every new open stratum is a torus bundle over an old open stratum, hence is aspherical by the homotopy exact sequence. Its normal-meridian lattice maps by a linearly independent set of primitive fan vectors into the old lattice, hence still injects into $\pi_1(U)$. This proves goodness at every step, and also proves the same preservation for the further stratal blow-ups below.

By D.8 every component $E_j$ of $E$ has log discrepancy zero over $(X,D)$. For a smooth SNC pair, the divisorial valuations of log discrepancy zero are exactly those produced by successive blow-ups of strata. One way to see the equality criterion is to pull back the logarithmic volume form: a center lying properly inside an open stratum supplies at least one ordinary coordinate and gives strictly positive discrepancy. A sequence following the center of the valuation therefore has discrepancy zero only at corner centers. In a normal orthant this is precisely a primitive monomial valuation. This criterion is also stated after Definition~2.7 in Svaldi's opened text, by reduction to the SNC discrepancy calculation. A primitive integral ray can be extracted by smooth star subdivisions. Write its coordinates in a smooth cone as $(a_1,\ldots,a_r)$ and blow up the face spanned by the basis vectors with positive coefficients. In the new cone indexed by a least positive coefficient $a_i$, the new coordinates are $a_i$ on the sum ray and $a_j-a_i$ on the other positive basis vectors. Their sum strictly decreases while at least two coefficients are positive. The process terminates at a coordinate ray, whose final coefficient is one by primitivity. Finitely many such extractions give a good smooth projective $X'\to X$ on which every prime divisor of $Y$ occurs.

Choose a very ample line bundle $H$ on $Y$. Transport its complete linear system to $X'$ and remove the common divisorial part. Let $I$ be the remaining base ideal and $L$ its line bundle. Its cosupport lies in $D'=X'\setminus U$ and has codimension at least two. Let
\[
 p:Z=\operatorname{Nor}(\operatorname{Bl}_I X')\longrightarrow X',
 \qquad q:Z\longrightarrow Y
\]
be the normalized graph. Writing $I\mathcal O_Z=\mathcal O_Z(-F)$, we have $q^*H=p^*L-F$ and $\mathcal O_Z(-F)$ is $p$-ample. Consequently $q$ is finite on every fibre of $p$: on such a fibre, the pullback of $H$ is ample, and a positive-dimensional fibre of $q$ would contain a curve on which that line bundle has degree zero.

Let $A$ be a $p$-exceptional prime divisor, equivalently a Rees divisor. It is absent as a divisor of $X'$, and therefore absent as a divisor of $Y$ by our extraction. Thus $\dim q(A)\leq1$. Its center on $X'$ cannot be a point, since then the two-dimensional $A$ would be mapped finitely by $q$ into a curve or a point. Hence its center is a curve $B$, and $q(A)$ is a curve $R$.

We now justify the crucial assertion that $B$ is a stratum curve and that the valuation is monomial. Choose a general point $b\in B$ avoiding every other exceptional center that does not dominate $B$, and a smooth surface slice transverse to $B$. Resolve the graph by principalizing $I$ on a smooth model. On the slice the resolution is a succession of point blow-ups of a smooth surface. Its exceptional divisor has a tree as dual graph, and its components are rational curves. The strict transforms of the boundary branches through $b$ give either one or two marked arrows. Each branch gives one arrow because its smooth curve germ has a single strict transform meeting the exceptional fibre in a single point. There are at most two branches at the generic point of a curve in a smooth SNC threefold.

Retain the vertices corresponding to Rees components, together with the marked arrows, and let $T_0$ be their minimal connecting subtree. A Rees component maps finitely to a rational curve $R_i\subset Y$. Let $S_i$ be the open stratum containing its generic point. The normalization of $R_i$ has at least two points outside $S_i$, by the absence of maps $\mathbb A^1\to S_i$ proved in D.8 (zero missing points also gives such a map). The same holds on its source rational component, because a finite surjection has a nonempty fibre over each of those two distinct points.

Every such boundary point comes from a boundary component $E_j$ not containing $R_i$. Its inverse image is an effective Cartier divisor on the normal graph $Z$, and therefore has pure codimension one. Here is the precise purity argument needed. Any supporting prime exceptional over $X'$ is a Rees divisor. Its center is a curve by the preceding paragraph; since $b$ avoids every other such center, that curve must be $B$. A supporting prime which is not exceptional is the strict transform of a prime divisor of $X'$. The map is the identity on $U$, so that prime lies in $D'$; it meets the general fibre over $b$ only when it contains $B$. These are exactly the one or two marked branches. The chosen Rees component is not itself in this support, since $E_j$ does not contain its image $R_i$.

On the resolved transverse surface, every exceptional component not corresponding to a Rees divisor contracts to a point of $Z$: a component mapping onto a curve of the general fibre would sweep out an exceptional prime of $Z$ dominating $B$, and would itself be a Rees component. Its connected non-Rees subtrees likewise contract. Thus all arrows and Rees components in any one direction of the resolved tree meet the chosen component at one point after those intervening subtrees are collapsed. A branch containing no arrow and no Rees component cannot supply an extra boundary point, because it supplies no supporting prime of the Cartier pullback on $Z$. It follows that every Rees vertex has at least two incident directions in $T_0$.

A finite tree with a retained vertex and only one marked arrow has an unmarked leaf. This contradicts the preceding assertion. Thus there are two arrows, so $B$ is an irreducible component of the intersection of two boundary divisors. Moreover $T_0$ is the path between the two arrows: a branch off this path would end at a Rees leaf. In the resolution of two transverse smooth curve germs, the vertices on that path are exactly those obtained by successive blow-ups of corner intersection points. Induction verifies this: a corner blow-up subdivides the path, while a blow-up at any other point adds a branch off it. Consequently every Rees valuation is a monomial valuation in the two boundary parameters. This also treats a Rees divisor whose generic fibre initially has several geometric components: the path is linearly ordered by its two labelled boundary arrows, and each corner component has its distinct primitive pair of weights. Monodromy preserves the two arrows and those weights, so it cannot permute different path positions. The valuation is therefore the indicated primitive monomial valuation over the function field of $B$.

For every $m\geq1$, the integral closure of $I^m$ is the intersection of the valuation ideals $v(f)\geq m v(I)$ over these finitely many Rees valuations. We use the existence and normalized-blow-up construction in Huneke--Swanson, Definition~10.1.1 and Theorem~10.2.2 \cite{HunekeSwanson2006}, both opened. Each inequality is monomial in boundary parameters, so each integral closure, and in particular $\overline I$, is a toroidal ideal sheaf. This includes points where several center curves meet. In a triple-crossing chart $x,y,z$, a Rees valuation over $x=y=0$ has weights $(a,b,0)$, and its valuation ideal is generated by the monomials $x^i y^j z^k$ with $ai+bj\geq m v(I)$. Valuations over $x=z=0$ or $y=z=0$ give the other two kinds of monomial inequality. At the special point the local integral closure is the intersection of all these valuation ideals. A finite intersection of monomial ideals is monomial: its generators are the least common multiples of choices of generators. The same argument applies to every subcollection of the stratum curves. An additional isolated integral-closure condition would require a Rees valuation centered at that point, already excluded by finiteness on $p$-fibres. Thus the normalized blow-up is toroidal everywhere, including the special triple points.

Principalize this monomial ideal by smooth stratal blow-ups. Canonical principalization, in the form recalled in AKMW \S1.2.3 \cite{AKMW2002}, is invariant under local torus actions and \'etale charts. Applied to a monomial ideal its smooth centers are unions of torus orbits closed in the chart, hence closures of strata; the local procedures agree on overlaps. We obtain a good $X_s\to X'$ and a morphism $X_s\to Y$. D.8 makes the latter toroidal. The corrected toroidal weak factorization theorem, AKMW Theorem~2.7.1 and Matsuki's erratum \S2 \cite{Matsuki2000}, factors it into smooth stratal blow-ups and blow-downs.

We use its relative morphism case. For the projective toroidal morphism $X_s\to Y$, the cobordism has a fan map to the fan of $Y$; its subdivisions preserve this map, and every elementary quotient fan still maps to that endpoint fan. These are the explicit fan maps in the proof of AKMW Theorem~2.7.1. Thus every intermediate variety is projective over $Y$ and its fan subdivides the fan of $Y$. This relative assertion proves goodness through the blow-down steps: every intermediate open stratum is a torus bundle over an open stratum of $Y$, and its meridian lattice injects, by the independent primitive generators of its cone, into the normal lattice of that endpoint stratum. Asphericity and meridian injectivity follow directly from goodness of $Y$, independently of the direction of the adjacent blow-up arrow. This gives the entire stated zigzag.

The proof uses only weak factorization. The algorithm for strong toroidal factorization in the earlier Morelli/AMR citation was withdrawn in the opened erratum and is not used.
\end{proof}

\begin{result}{D.10}{Proved}{All one-step stratal modifications of $\overline M_{0,6}$}
There are exactly $210$ nontrivial labelled smooth stratal blow-ups of $\overline M_{0,6}$: $105$ along boundary curves and $105$ at boundary points. Under permutation of the six labels they form four orbits, of sizes $45,60,15,90$. There is no nontrivial one-step stratal blow-down from the standard pair to another smooth SNC pair. The associated maps of conical complexes are the star subdivisions at the chosen two- or three-dimensional cones, with new ray the sum of their primitive generators. Theorem D.9 consequently applies to every good smooth projective completion of $M_{0,6}$.
\end{result}
\begin{proof}
A stable six-legged tree has at most three edges. The exact enumeration in \path{code/d_boundary_checks.py}, using the independently generated graph list of Part~A, gives $25,105,105$ strata of codimensions one, two, three. Blowing up a divisor is the identity. For two-edge trees the end/middle leg distributions are $(2,2,2)$ and $(2,1,3)$, with $45$ and $60$ labellings. For three-edge trees the internal graph is either a three-armed star, with two legs at each endpoint, or a path with endpoint leg counts two and interior leg counts one. Their counts are $15$ and $90$. Smoothness of the corresponding star subdivisions follows because each replacement maximal-cone matrix has determinant $1$; all these determinants are checked exactly by the program.

A smooth stratal blow-down in dimension three must contract a boundary divisor that is either a projective plane over a point or a projective-line bundle over a smooth curve. A divisor of type $2|4$ is $\overline M_{0,5}$, a rational surface with second Betti number five. It is neither $\mathbb P^2$ nor a projective-line bundle over a curve: rationality forces that curve to have genus zero, and such a bundle has second Betti number two. A divisor of type $3|3$ is $\mathbb P^1\times\mathbb P^1$. Any contraction to a curve uses a ruling, but its boundary contains three sections of that ruling. Their images would be three distinct boundary branches containing the entire center curve. Three SNC branches in a smooth threefold meet in codimension three, so this target would fail the SNC condition. These exhaust the twenty-five divisors. Finally, at each stratum of the standard moduli compactification the cut-surface Teichm\"uller groups are aspherical and independent commuting Dehn twists inject; hence the pair is good and D.9 applies.
\end{proof}
